\section{Two-snapshot ancestry}
\label{sec:problem}

\paragraph{Task.}
We observe a population at times $a$ and $b{=}a{+}H$.
Each snapshot provides cell detections $\mathcal{A}{=}\{\mathbf{x}_i\}_{i=1}^{N_a}$ and $\mathcal{B}{=}\{\mathbf{y}_j\}_{j=1}^{N_b}$ (centroid, and optionally major/minor axis and orientation) and, when available, the image.
For every target $j$ we predict its ancestor $\pi(j)\in\mathcal{A}\cup\{\varnothing\}$, the cell at time $a$ whose lineage contains $j$; $\varnothing$ denotes cells that entered the field of view.
No intermediate frame and no identity information is available at inference.
Ground-truth ancestry comes from full-resolution tracking and is used only for training and evaluation.

\paragraph{Metrics.}
Accuracy is the fraction of evaluable targets (inside the field of interest, with a visible ancestor) assigned to the correct ancestor, averaged over windows (one snapshot pair of one movie); weighting by cells instead counts dense windows more and gives lower numbers (Supp.~E).
At the family level we compare biological readouts of the predicted and true families (size control; family-size spread in Supp.~H).
We report results separately on \emph{dense} windows ($N_a{>}200$), where ambiguity is highest.

\paragraph{Data.}
\cref{tab:datasets} summarises the datasets.
\emph{OIAM}~\cite{seiffarth2024oiam}: five \emph{C.~glutamicum} microcolony movies (M0--M4) with over 1.4M annotated instances; we train in-domain models on M0--M2 (or four movies in leave-one-movie-out) and evaluate at $H\in\{33,66,132\}$ minutes.
\emph{FAST}~\cite{meacock2023fast}: seven \emph{E.~coli} microcolonies with manually corrected lineages.
\emph{recA} and \emph{rpsM}: nine \emph{E.~coli} microcolonies each from the raw data of~\cite{vanvliet2018spatially} (also FAST's source), with other reporters, dates and colonies and uncurated Schnitzcells lineages (audited in Supp.~E).
\emph{LFCT}~\cite{gachloo2026lfct}: one field of motile MDA-MB-231 breast cancer cells; \emph{T98G}~\cite{chang2026tui}: one field of T98G glioblastoma cells under electrotaxis.
The last two are sparse and motile, a regime deliberately unlike the simulator.
External datasets are never used for training.
The window lists, the \rodsim simulator, code and trained models will be released.

\begin{table}[t]
\centering\small
\setlength{\tabcolsep}{3pt}
\resizebox{\linewidth}{!}{%
\begin{tabular}{lllrrr}
\toprule
Dataset & Organism & Gaps (min) & Seq. & Windows & max $N_a$ \\
\midrule
OIAM (M3/M4) & \emph{C. glutamicum} & 33/66/132 & 2 & 109 & 1359 \\
FAST & \emph{E. coli} & 35/65/130 & 7 & 1409 & 124 \\
recA & \emph{E. coli} & 30/60/120 & 9 & 345 & 108 \\
rpsM & \emph{E. coli} & 30/60/120 & 9 & 900 & 113 \\
LFCT & MDA-MB-231 & 30/60/120 & 1 & 1141 & 13 \\
T98G & T98G & 30/70/130 & 1 & 88 & 74 \\
\bottomrule
\end{tabular}}
\caption{\textbf{Evaluation datasets.} Windows pool all gaps.}
\label{tab:datasets}
\end{table}
